%%%%%%%%%%  scan idx 202  |  folio 189  %%%%%%%%%%
\emph{Preuve.} Comme $f$ est quasi-compact et quasi-séparé (EGA IV 1.2.3 et 1.2.4), $f_*$ transforme faisceaux quasi-cohérents en faisceaux quasi-cohérents (III 1.4.10 et IV 1.7.21). Les foncteurs $Qf_*$ et $f_*Q$ de $\Mod(X)$ dans $\Qcoh(Y)$ sont tous deux adjoints à droite du foncteur $f^*:\Qcoh(Y)\longrightarrow\Mod(X)$, donc canoniquement isomorphes.

\textbf{Proposition 3.5.} Si $S$ est noethérien, ou quasi-compact et séparé, le foncteur
\begin{equation}\tag{3.0}
 \phi:\mathrm D^+(\Qcoh(S))\longrightarrow \mathrm D^+(S)
\end{equation}
est pleinement fidèle et a pour image essentielle la sous-catégorie pleine $\mathrm D^+(S)_{\qcoh}$ de $\mathrm D^+(S)$ formée des complexes à cohomologie quasi-cohérente.

\emph{Preuve.} On peut se borner à prouver l'assertion relative au cas où $S$ est séparé, le cas où $S$ est noethérien étant bien connu ([H] II 7.19). Il revient évidemment au même de prouver les assertions suivantes :

\begin{enumerate}[label={(3.5.\arabic*)~:}]
\item pour $E\in\Ob\mathrm D^+(\Qcoh(S))$, la flèche d'adjonction
\[
 E\longrightarrow \mathbb RQ(E)
\]
(déduite de (3.3 b) (ii)) ) est un isomorphisme ;
\item pour $E\in\Ob\mathrm D^+(S)_{\qcoh}$, la flèche d'adjonction
\[
 \mathbb RQ(E)\longrightarrow E
\]
(déduite de (3.3 b)(ii)) ) est un isomorphisme.
\end{enumerate}

\emph{Preuve de} (3.5.1). Par un dévissage immédiat\footnote{Utilisant par exemple la suite spectrale $R^pQ(H^q(E))\Longrightarrow R^{*}Q(E)$.}, on est ramené à prouver que les faisceaux quasi-cohérents sont acycliques pour le cohéreur, (i.e. que, pour $E\in\Ob\Qcoh(S)$, on a $R^iQ(E)=0$ pour $i>0$), et que d'autre part, pour $E\in\Ob\Qcoh(S)$, la flèche d'adjonction $E\longrightarrow QE$ est un isomorphisme. Si $S$ est affine, cela résulte de la formule $Q=\widetilde{\Gamma}_S$ (cf. preuve de (3.2)). Dans le cas général, recouvrons $S$ par un nombre fini d'ouverts affines $U_i$. Comme $S$ est séparé, les inclusions $f_i:U_i\longrightarrow S$ sont des morphismes affines, 
